\section{TCSA: Temporal-Capability Safety Architecture}
\label{sec:tcsa}

The preceding layers defend against attacks that are individually
identifiable---a malicious pattern, a forbidden capability, a suspicious
behavioral signature. But a critical class of attacks evades every such layer:
\emph{within-authority chaining}, where each individual action is fully
legitimate yet the composition constitutes a security violation.

\paragraph{Motivating example.} Consider the following four-step tool chain
executed by an LLM agent:

\begin{enumerate}[leftmargin=*]
\item \texttt{file\_read(".env")} --- reads the application's environment file.
      \emph{Legitimate}: the agent has file-read capability.
\item \texttt{parse(output, "credentials")} --- extracts database credentials
      from the file contents. \emph{Legitimate}: parsing is a core LLM
      capability.
\item \texttt{compose\_email(to="external@evil.com", body=credentials)} ---
      drafts an email containing the extracted credentials. \emph{Legitimate}:
      the agent has email-composition capability.
\item \texttt{send\_email()} --- transmits the composed email.
      \emph{Legitimate}: the agent has email-send capability.
\end{enumerate}

Every action is individually authorized. No single layer catches the
violation. Yet the chain constitutes \textbf{data exfiltration}---sensitive
credentials have been read, composed into a message, and transmitted to an
external attacker-controlled address. The attack is invisible to any defense
that examines actions in isolation.

\tcsa addresses this gap through three sub-primitives that collectively
monitor temporal ordering, data-flow capabilities, and predictive trajectory:

\begin{itemize}[leftmargin=*]
\item \textbf{\tsa} (Temporal Safety Automata) --- adapted from runtime
      verification; monitors temporal ordering of tool calls against LTL
      safety properties.
\item \textbf{\cafl} (Capability-Attenuating Flow Labels) --- a genuinely new
      primitive; tracks data capabilities through LLM transformations under
      worst-case taint assumptions.
\item \textbf{\gps} (Goal Predictability Score) --- a genuinely new primitive;
      provides predictive early warning by scoring how many reachable future
      states lead to safety violations.
\end{itemize}

\subsection{TSA: Temporal Safety Automata}
\label{sec:tcsa:tsa}

\tsa adapts techniques from runtime verification---specifically the
monitor-automata approach of Havelund and Ro\c{s}u~\cite{havelund2004} and
Pnueli's Linear Temporal Logic (LTL)~\cite{pnueli1977}---to the setting of
LLM tool chains. While runtime verification has been extensively studied for
software systems, it has never been applied to monitor sequences of LLM tool
invocations. The adaptation is non-trivial: unlike deterministic programs,
LLM agents select actions stochastically, and the space of possible tool
chains is combinatorially large.

\paragraph{Design-time compilation.} Safety properties are expressed as LTL
formulae over the alphabet of tool-call events. At design time, each formula
is compiled into a deterministic monitor automaton via standard
LTL-to-B\"uchi-to-monitor construction~\cite{havelund2004}. The resulting
automata are minimal and fixed-size, requiring no runtime compilation.

\paragraph{Example safety properties.} The following three properties
illustrate the expressiveness of the \tsa formalism:

\begin{align}
  \mathbf{P_1}\!: \quad & \Box\bigl(\mathit{read\_sensitive}
    \rightarrow \Box\,\neg\,\mathit{send\_external}\bigr)
    \label{eq:p1} \\
  \mathbf{P_2}\!: \quad & \neg\Diamond\bigl(\mathit{read\_credentials}
    \wedge \Diamond\,\mathit{send\_external}\bigr)
    \label{eq:p2} \\
  \mathbf{P_3}\!: \quad & \Box\bigl(\mathit{privilege\_change}
    \rightarrow \bigcirc\,\mathit{approval\_received}\bigr)
    \label{eq:p3}
\end{align}

Property~$\mathbf{P_1}$ states that after reading sensitive data, the agent
may \emph{never} subsequently send data to an external channel---a permanent
prohibition. Property~$\mathbf{P_2}$ forbids any execution trace in which
credential access is eventually followed by external transmission---directly
encoding the exfiltration scenario from the motivating example.
Property~$\mathbf{P_3}$ requires that every privilege escalation be
\emph{immediately} followed by an approval event, enforcing a mandatory
human-in-the-loop checkpoint.

\paragraph{Abstract security state.} To make monitor evaluation tractable, we
define an \emph{abstract security state} as a 16-bit vector encoding the
security-relevant predicates of the current execution. The 16~bits encode
predicates including:

\begin{itemize}[leftmargin=*]
\item \texttt{sensitive\_data\_accessed} --- set when any tool reads data
      classified at \textsc{internal} or above.
\item \texttt{credentials\_accessed} --- set when authentication material
      (API keys, passwords, tokens) is read.
\item \texttt{external\_channel\_opened} --- set when a tool establishes a
      connection to an external endpoint.
\item \texttt{privilege\_elevated} --- set when the agent's capability set
      expands beyond its initial grant.
\item \texttt{write\_performed} --- set when any destructive or
      state-modifying tool call executes.
\item \texttt{approval\_pending} --- set when a privilege change awaits
      operator confirmation.
\end{itemize}

A 16-bit vector yields $2^{16} = 65{,}536$ possible abstract states. This
state space is small enough for \emph{complete enumeration} at design time,
enabling pre-computation of all monitor transitions and, critically, enabling
the \gps predictive analysis described in~\S\ref{sec:tcsa:gps}.

\paragraph{Runtime cost.} At runtime, each tool call triggers a single
state-vector update (bitwise OR of the relevant predicate bits) followed by a
table lookup in the pre-compiled monitor transition function. Both operations
are $O(1)$ with respect to chain length, tool count, and property count.
Monitor evaluation adds negligible overhead---under 1\,$\mu$s per tool
call---making \tsa suitable for latency-sensitive deployments.

\subsection{CAFL: Capability-Attenuating Flow Labels}
\label{sec:tcsa:cafl}

\cafl is a genuinely new primitive with no direct precedent in the literature.
Classical information flow control (IFC)~\cite{denning1976} tracks how data
propagates through deterministic programs, where the transformation applied to
each datum is known and analyzable. This assumption breaks catastrophically
for LLM-based agents: the model can perform \emph{any} transformation on its
input, including arbitrary summarization, recombination, and rephrasing. No
static analysis can determine what an LLM will do with a given input.

\paragraph{Key insight: worst-case taint.} \cafl resolves this by adopting
a \emph{worst-case taint} assumption: if data carrying capability label~$C$
enters an LLM, \textbf{all output from that LLM invocation is tainted with
the restrictions of~$C$}. This is deliberately conservative---it may
over-approximate (flagging safe chains as dangerous) but will
\emph{never} under-approximate (allowing dangerous chains through). Soundness
is preserved at the cost of completeness, which is the correct trade-off for
a security mechanism.

\paragraph{Capability labels.} Every data object in the system carries a
\emph{capability label}---a set of permitted downstream operations. The
critical invariant is that capabilities may only \emph{decrease} as data flows
through the processing chain; they may never increase. This monotonic
attenuation ensures that sensitive data progressively loses its ability to
reach dangerous sinks:

\begin{itemize}[leftmargin=*]
\item $\texttt{file\_read(".env")} \rightarrow \text{output:}~
      \{\textsc{process}, \textsc{display}\}$\\
      Sensitive configuration data may be processed internally and displayed
      to the operator, but \emph{cannot} be exported.
\item $\texttt{file\_read("public.md")} \rightarrow \text{output:}~
      \{\textsc{process}, \textsc{display}, \textsc{export}\}$\\
      Public data carries the full capability set, including export.
\item $\texttt{email\_send()} ~\text{requires input:}~
      \{\textsc{export}\} \in C$\\
      The email tool checks that its input data carries the \textsc{export}
      capability before transmission.
\end{itemize}

\paragraph{Blocking the exfiltration chain.} Revisiting the motivating
example: the chain \texttt{.env} $\rightarrow$ LLM $\rightarrow$
\texttt{email\_send} is blocked because \texttt{file\_read(".env")} produces
output labeled $\{\textsc{process}, \textsc{display}\}$. When this data
enters the LLM, worst-case taint applies: the LLM's output inherits the
\emph{intersection} of input capabilities, yielding at most
$\{\textsc{process}, \textsc{display}\}$. The subsequent
\texttt{email\_send()} requires $\{\textsc{export}\} \in C$, which is absent.
The chain is rejected before the email is ever composed.

\paragraph{Membrane pattern.} \cafl enforces capability attenuation at
\emph{trust boundary crossings}---points where data moves between security
domains (e.g., from internal processing to external communication, or from
the agent's workspace to a third-party API). Each boundary acts as a
\emph{membrane} that filters capabilities: data may cross only if it carries
the capabilities required by the destination domain. This membrane pattern
provides defense in depth, as even a compromised internal component cannot
elevate the capabilities of data that has already been attenuated.

\subsection{GPS: Goal Predictability Score}
\label{sec:tcsa:gps}

\gps is a genuinely new primitive that provides \emph{predictive} defense:
it catches tool chains that are \emph{heading toward} a safety violation
before the violation actually occurs. Where \tsa detects violations at the
moment they happen (the monitor enters a rejecting state), \gps provides
early warning by analyzing the \emph{trajectory} of the current execution
through the abstract state space.

\paragraph{Definition.} Given the current abstract security state
$s \in \{0,1\}^{16}$, define $\mathcal{N}(s)$ as the set of states reachable
from~$s$ via a single tool call, and $\mathcal{D}(s) \subseteq \mathcal{N}(s)$
as the subset of those states that either (a)~directly violate a \tsa safety
property or (b)~have no safe successor under any tool call. The Goal
Predictability Score is:

\begin{equation}
  \gps(s) = \frac{|\mathcal{D}(s)|}{|\mathcal{N}(s)|}
  \label{eq:gps}
\end{equation}

\noindent A score of $\gps(s) = 0$ indicates that no single-step successor
leads to danger; $\gps(s) = 1$ indicates that \emph{every} possible next
action leads to a safety violation.

\paragraph{Tractability.} The key observation enabling \gps is that the
abstract state space is only $2^{16} = 65{,}536$ states. This is small enough
for \emph{complete enumeration}: at design time, we pre-compute
$\mathcal{D}(s)$ and $\mathcal{N}(s)$ for every state~$s$, storing the
resulting \gps values in a lookup table. At runtime, computing \gps requires
a single table lookup---the same $O(1)$ cost as the \tsa monitor transition.

\paragraph{Threshold policy.} We define a two-threshold policy based on
the \gps value:

\begin{itemize}[leftmargin=*]
\item $\gps(s) > 0.7$: \textbf{Warning.} The system alerts the operator that
      the current tool chain is approaching a region of the state space where
      most continuations lead to safety violations. The agent may continue,
      but its actions are logged at elevated verbosity.
\item $\gps(s) > 0.9$: \textbf{Block.} The system pre-emptively halts the
      tool chain. Fewer than 10\% of possible continuations are safe,
      indicating that the agent has maneuvered into a near-certain violation
      trajectory regardless of its next action.
\end{itemize}

\paragraph{Complementarity with \tsa.} \gps and \tsa are complementary:
\tsa provides \emph{reactive} enforcement (blocking the exact moment a
property is violated), while \gps provides \emph{proactive} defense (warning
or blocking before the violation occurs). Together, they implement a
defense-in-depth strategy over the temporal dimension: \gps catches attacks
that are still developing, and \tsa catches any that slip past the predictive
threshold.

\subsection{Comparison with Existing Approaches}
\label{sec:tcsa:comparison}

The closest existing work to \tcsa is CrossToolGuard, which analyzes pairs of
consecutive tool calls for dangerous combinations. Table~\ref{tab:tcsa-comparison}
summarizes the key differences. \tcsa subsumes pair-based analysis while
adding temporal ordering, data-flow tracking, and predictive capabilities
that no existing system provides.

\begin{table}[t]
\centering
\small
\caption{Comparison of \tcsa with pair-based tool-chain analysis
(CrossToolGuard). \tcsa provides strictly stronger guarantees across every
dimension.}
\label{tab:tcsa-comparison}
\begin{tabular}{@{}lll@{}}
\toprule
Aspect & CrossToolGuard & \tcsa \\
\midrule
Chain length    & Pairs only      & Arbitrary length \\
Temporal order  & No              & Yes (LTL via \tsa) \\
Data flow       & No              & Yes (\cafl labels) \\
Predictive      & No              & Yes (\gps score) \\
Runtime cost    & $O(N^2)$ pairs  & $O(1)$ per call \\
\bottomrule
\end{tabular}
\end{table}

The $O(1)$-per-call runtime of \tcsa deserves emphasis. CrossToolGuard must
evaluate every pair of tools in the current chain, yielding $O(N^2)$ cost for
a chain of length~$N$. \tcsa maintains a single 16-bit state vector and
performs one table lookup per tool call, regardless of chain length or the
number of safety properties being monitored. This constant-time guarantee
makes \tcsa viable even for long-running agentic sessions with hundreds of
sequential tool invocations.
